\section{Payroll Bridge and Exposure-Construction Diagnostics}\label{app:bcc_architecture}

\subsection{BCC version and grouping bridge}

The comparison uses the August 12, 2026 version of Brynjolfsson--Chandar--Chen and records the accessed paper, dashboard, and workbook in the evidence ledger.  Their reported 19-percent young-worker ``kept pace'' shortfall is a stock-growth counterfactual in ADP data, not a regression coefficient, employment probability, or hiring rate.  Their public dashboard forms equal-occupation exposure groups; later regressions employ other weights and employer controls.  Exact dashboard memberships and the treatment of ties are not released.

The CPS bridge therefore applies an equal-occupation, tie-preserving top-two-versus-bottom-three approximation on the present support and labels it as such.  The pooled estimate is $-0.0720$, with interval $[-0.1223,-0.0218]$.  The SOC2-by-calendar-month estimate is $-0.0168$, with interval $[-0.0808,0.0473]$.  Their paired movement is 0.0553, with interval $[-0.0013,0.1118]$.  The intermediate SOC2-by-post specification moves by 0.0573 relative to the pooled model, with interval $[0.0005,0.1141]$; this post-outcome bridge contrast is detected narrowly, whereas the stronger family-month comparison is not.  The archive provides group summaries and a membership file; it cannot provide concordance with undisclosed assignments.  Differences in occupation universe, age comparison, denominator, endpoint, firm controls, and employer-matched flows remain.

\begin{table}[!htbp]
\centering
\caption{Post-outcome exploratory BCC-style grouping bridge}\label{tab:appendix_bcc_bridge}
\scriptsize
\setlength{\tabcolsep}{3pt}
\resizebox{\textwidth}{!}{%
\begin{tabular}{llrrr}
\toprule
Grouping rule & Conditioning structure & Coefficient [95\% CI] & Change from baseline [95\% CI] & Paired $\mathrm{MDE}_{80}$ \\
\midrule
\multicolumn{5}{l}{\textit{Panel A. Public-dashboard equal-occupation approximation}} \\
Equal occupation & Occupation and calendar-month FE & $-0.072\;[-0.122,-0.022]$ & --- & --- \\
Equal occupation & SOC2 by young by post & $-0.015\;[-0.080,0.050]$ & $0.057\;[-0.001,0.115]$ & 0.083 \\
Equal occupation & SOC2 by young by calendar month & $-0.017\;[-0.081,0.047]$ & $0.055\;[-0.002,0.113]$ & 0.082 \\
\addlinespace
\multicolumn{5}{l}{\textit{Panel B. Historical employment-weighted approximation}} \\
Employment weighted & Occupation and calendar-month FE & $-0.075\;[-0.125,-0.025]$ & --- & --- \\
Employment weighted & SOC2 by young by post & $-0.027\;[-0.095,0.041]$ & $0.048\;[-0.011,0.108]$ & 0.086 \\
Employment weighted & SOC2 by young by calendar month & $-0.027\;[-0.093,0.040]$ & $0.048\;[-0.011,0.107]$ & 0.086 \\
\bottomrule
\end{tabular}}
\tabnote{All rows are post-outcome exploratory CPS bridges, not replications of BCC's proprietary ADP analysis.  Each row compares GPT-4-beta Q4--Q5 with Q1--Q3 on the same 468 occupations, includes the standardized historical Webb-software control, and uses 9,999 common occupation-level Rademacher multipliers.  Paired changes compare each SOC2-conditioned row with its own unconditioned baseline.  Panel A follows the public dashboard's equal-occupation description as closely as public information permits.  Panel B uses earlier employment-weighted cuts formed from 108 months of young-plus-older CPS stock, including postperiod months, and is only a grouping-rule diagnostic.  Exact BCC occupation membership, universe, cutoff and tie algorithm are not public, so concordance is unverified.  The BCC and CPS objects also differ in ages, denominator, endpoint, employer controls, and observed worker--firm flows.  Null-containing paired intervals are nondetections, not equivalence.}
\end{table}


\subsection{Annual ACS reconstruction}

The annual extension uses ordinary one-year ACS PUMS for 2017--2019 and
2021--2024; no standard 2020 one-year file exists, and no experimental 2020
estimate is spliced into the panel.  Eighty replicate weights per endpoint year
form the successive-difference-replication intervals in
Table~\ref{tab:appendix_acs}.  The BCC-labelled rows remain public analogues:
complete BCC occupation membership is unavailable, so they cannot be exact
replications.

\begin{table}[!htbp]
\centering
\caption{Independent annual ACS reconstructions}\label{tab:appendix_acs}
\scriptsize
\setlength{\tabcolsep}{2pt}
\resizebox{\textwidth}{!}{%
\begin{tabular}{llllrrr}
\toprule
Definition & Population & Calendar/model & Estimate [ACS CI] & Family-year [ACS CI] & Family-year occupation CI & Family-year family CI \\
\midrule
\multicolumn{7}{l}{\textit{Panel A. 2022--2024 Q5--Q1 growth-factor benchmark, ages 22--25}} \\
BCC published & All employed & 2022--2024 & $-0.022\;[-0.055,0.011]$ & --- & --- & --- \\
ACS analogue, equal occupation & All employed & 2022--2024 & $-0.0859\;[-0.1238,-0.0480]$ & --- & --- & --- \\
Aligned annual CPS analogue & All employed & 2022--2024 & $-0.0346\;[-0.1455,0.0763]$ & --- & --- & --- \\
BCC published & Full-time civilian wage/salary & 2022--2024 & $-0.019\;[-0.059,0.020]$ & --- & --- & --- \\
ACS analogue, equal occupation & Full-time civilian wage/salary & 2022--2024 & $-0.0834\;[-0.1272,-0.0396]$ & --- & --- & --- \\
Aligned annual CPS analogue & Full-time civilian wage/salary & 2022--2024 & $-0.0306\;[-0.1366,0.0755]$ & --- & --- & --- \\
\addlinespace
\multicolumn{7}{l}{\textit{Panel B. Annual young-relative Q5-by-post regressions, ages 22--25 versus 26--65}} \\
Fixed primary groups & All employed & Full & $-0.1089\;[-0.1310,-0.0869]$ & $-0.0474\;[-0.0885,-0.0063]$ & $[-0.1316,0.0368]$ & $[-0.1242,0.0294]$ \\
Fixed primary groups & All employed & Nonreuse & $-0.0951\;[-0.1204,-0.0698]$ & $-0.0220\;[-0.0705,0.0264]$ & $[-0.1121,0.0680]$ & $[-0.1013,0.0572]$ \\
Fixed primary groups & Full-time civilian wage/salary & Full & $-0.1196\;[-0.1456,-0.0936]$ & $-0.0691\;[-0.1193,-0.0188]$ & $[-0.1554,0.0173]$ & $[-0.1777,0.0395]$ \\
Fixed primary groups & Full-time civilian wage/salary & Nonreuse & $-0.1063\;[-0.1361,-0.0764]$ & $-0.0467\;[-0.1057,0.0122]$ & $[-0.1317,0.0383]$ & $[-0.1480,0.0546]$ \\
Broader equal-occupation groups & All employed & Full & $-0.0830\;[-0.1053,-0.0607]$ & $-0.0151\;[-0.0559,0.0257]$ & $[-0.1000,0.0698]$ & $[-0.0942,0.0639]$ \\
Broader equal-occupation groups & All employed & Nonreuse & $-0.0723\;[-0.0979,-0.0467]$ & $-0.0056\;[-0.0535,0.0424]$ & $[-0.0918,0.0806]$ & $[-0.0927,0.0815]$ \\
Broader equal-occupation groups & Full-time civilian wage/salary & Full & $-0.0933\;[-0.1203,-0.0664]$ & $-0.0292\;[-0.0790,0.0205]$ & $[-0.1112,0.0527]$ & $[-0.1307,0.0722]$ \\
Broader equal-occupation groups & Full-time civilian wage/salary & Nonreuse & $-0.0828\;[-0.1137,-0.0520]$ & $-0.0291\;[-0.0875,0.0293]$ & $[-0.1118,0.0535]$ & $[-0.1294,0.0712]$ \\
\bottomrule
\end{tabular}}
\tabnote{Panel A reports growth-factor differences without an older comparison group.  The public analogues do not recover BCC's undisclosed complete occupation membership.  Panel B reports log coefficients.  Full uses 2017--2019 and 2021 as preperiod years and 2023--2024 as postperiod years; nonreuse uses 2017 and 2021 as preperiod years.  Standard 2020 one-year ACS PUMS are omitted, and 2022 is the transition year.  ACS intervals use successive-difference replication with 80 replicate weights per perturbed endpoint year and a year-block independence approximation.  Occupation and family intervals instead assign shocks at those cluster levels; the procedures answer different stochastic questions and are not combined.  All extensions are post-outcome descriptive analyses.}
\end{table}


For the benchmark, changing from all employed to full-time civilian
wage-and-salary workers is not detected under common ACS replicate
perturbations: the paired differences are $-0.0025$
($[-0.0266,0.0216]$) for the equal-occupation BCC analogue and $0.0056$
($[-0.0174,0.0285]$) for fixed primary groups.  The fixed-primary versus BCC-analogue
grouping difference is likewise not detected among all employed workers
($0.0090$, $[-0.0045,0.0224]$).  Among full-time civilian wage-and-salary
workers that grouping difference is narrowly detected ($0.0171$,
$[0.0002,0.0340]$), which reinforces rather than resolves the missing-membership
limitation.

For the annual young-relative regressions, full-calendar family-year estimates
retain negative ACS sampling intervals under the fixed primary exposure groups,
while the broader equal-occupation grouping and every nonreuse-calendar
family-year interval contain zero.  The paired ACS sampling intervals for the
pooled-minus-family-year movements exclude zero in all eight reported
definition--population--calendar combinations.  This is not universal across
stochastic targets: supplementary occupation- and family-multiplier intervals
are wider, and every family-year coefficient reported in
Table~\ref{tab:appendix_acs} includes zero under those shock assignments.  The
annual exercise thus shows that large respondent counts sharpen person-sampling
uncertainty but do not create within-family occupational comparisons or
identify an AI effect.

\subsection{Equal-occupation estimation objective}

The public-dashboard grouping bridge above changes the grouping rule.  A separate sensitivity changes the regression objective itself: each occupation receives total weight one, while \texttt{WTFINL} still determines the young and older shares within occupation--month cells.  Under this objective the pooled Q5-by-post coefficient is $-0.0833$, compared with $-0.1321$ under employment-stock weighting.  Their paired difference is $0.0488$ ($[-0.0604,0.1581]$).  The corresponding family-month estimates are $-0.0047$ and $-0.0217$, with a paired difference of $0.0170$ ($[-0.1587,0.1927]$).  The design therefore does not detect an objective-driven difference, but the wider equal-occupation intervals and null-containing paired intervals do not establish equivalence.

\begin{table}[!htbp]
\centering
\caption{Equal-occupation objective sensitivity}\label{tab:appendix_equal_objective}
\small
\setlength{\tabcolsep}{4pt}
\resizebox{\textwidth}{!}{%
\begin{tabular}{llrr}
\toprule
Conditioning structure & Objective & Q5 $\times$ post [95\% CI] & Equal occupation $-$ stock weighted [95\% CI] \\
\midrule
Occupation and calendar-month FE & Employment-stock weighted & $-0.132\;[-0.220,-0.045]$ & --- \\
Occupation and calendar-month FE & Equal occupation & $-0.083\;[-0.206,0.040]$ & $0.049\;[-0.060,0.158]$ \\
SOC2 family by calendar month & Employment-stock weighted & $-0.022\;[-0.159,0.115]$ & --- \\
SOC2 family by calendar month & Equal occupation & $-0.005\;[-0.203,0.193]$ & $0.017\;[-0.159,0.193]$ \\
\bottomrule
\end{tabular}}
\tabnote{This sensitivity changes the estimation objective: every occupation receives total weight one, while \texttt{WTFINL} continues to determine young and older shares within occupation--month cells.  Treatment membership, 468-occupation support, calendar, and nuisance structures are unchanged.  Intervals use 9,999 common occupation-level Rademacher score multipliers; paired differences are formed before applying the common draws.  A paired interval containing zero means that this design does not detect a difference, not that the objectives are economically equivalent.}
\end{table}


\subsection{Task-capability weighting}

Write the two Eloundou primitives as $E_1$ (the E1 share) and $E_2$ (the E2 share).  The score family $E_1+\lambda E_2$ is evaluated on a fixed grid from zero to one under common preperiod weights and common bootstrap draws.  The implementation verifies numerically that $\lambda=0.5$ reproduces raw beta before grouping.  Any earlier mismatch between that row and the headline arose from a different support or treatment contract and is preserved only as implementation history.

Two versions are kept separate.  The native-rank analysis recomputes the score normalization and tie-preserving groups at each $\lambda$, so both scale and membership can change.  The fixed-membership analysis changes only the continuous score representation while holding the beta groups fixed.  For every grid point the result files report coefficient, paired difference from beta, interval, occupation coverage, tail turnover, and correlations with computer use and remotability.  A focused 408-occupation run adds both computer use and remotability on common support.  None of these paths measures realized deployment.

\subsection{Primitive and architecture comparisons}

The joint primitive model reports $E_1$ and $E_2$ in raw share units, in preperiod-standard-deviation units, their covariance, and an illustrative common change.  This is a two-regressor parameterization, not a causal decomposition into technologies.  Architecture comparisons likewise report both component coefficients and their paired difference on each declared support.  Webb-AI comparisons are shown first on identical support and then on broader native support so that score changes are not confused with population changes.  Coordinate rotations that do not change the fitted predictor are not interpreted as additional scientific results.

Raw AIOE is a standardized source-occupation index: one point is approximately one unweighted standard deviation across the 774 occupations in the source workbook, and the source interquartile difference is 1.879 points.  It is not a percentage.  Equal-administrative and source-employment-weighted mappings preserve that source scale through averaging, although their target-support distributions need not retain unit variance.  The direct-ability reconstruction has its own raw scale and is not numerically interchangeable.  Categorical AIOE estimates use ranks and tie-preserving groups; continuous comparisons must state whether their increment is a raw point or a current-support weighted standard deviation.

Table~\ref{tab:appendix_nonlambda_architecture} runs the five non-$\lambda$ definitions on the current 113-month outcome contract.  Native-support pooled coefficients range from $-0.0991$ for direct-ability AIOE to $-0.0145$ for the reversed OECD gap.  Every native-support alternative-minus-beta paired occupation interval contains zero; the OECD comparison nearly reaches the boundary, with interval $[-0.0004,0.2378]$.  On the common 444-occupation AIOE support, however, family-month differences between administrative/equal-component mapping and direct ability ($0.0688$, $[-0.0070,0.1447]$ under occupation clustering but $[0.0057,0.1320]$ under family clustering) and source-employment weighting ($0.0586$, $[0.0111,0.1061]$ under occupation clustering) show that implementation can matter for the conditioned result.  These are paired implementation sensitivities within AIOE, not independent replications or causal technology decompositions.

\begin{table}[!htbp]
\centering
\caption{Non-$\lambda$ exposure definitions on the current CPS contract}\label{tab:appendix_nonlambda_architecture}
\scriptsize
\setlength{\tabcolsep}{3pt}
\resizebox{\textwidth}{!}{%
\begin{tabular}{lrrrrr}
\toprule
Exposure definition & Occupations & Pooled coefficient [95\% CI] & Alternative $-$ beta [95\% CI] & Family-month coefficient [95\% CI] & Alternative $-$ beta [95\% CI] \\
\midrule
AIOE, administrative/equal component & 466 & $-0.072\;[-0.143,-0.001]$ & $0.064\;[-0.025,0.152]$ & $0.065\;[-0.050,0.181]$ & $0.095\;[-0.063,0.253]$ \\
AIOE, direct ability & 455 & $-0.099\;[-0.174,-0.024]$ & $0.027\;[-0.059,0.113]$ & $0.024\;[-0.084,0.132]$ & $0.011\;[-0.143,0.165]$ \\
AIOE, source-employment weighted & 456 & $-0.072\;[-0.145,0.000]$ & $0.068\;[-0.020,0.157]$ & $0.071\;[-0.044,0.186]$ & $0.114\;[-0.047,0.274]$ \\
Webb artificial intelligence & 443 & $-0.068\;[-0.160,0.023]$ & $0.067\;[-0.036,0.171]$ & $-0.079\;[-0.177,0.019]$ & $-0.036\;[-0.203,0.131]$ \\
OECD reversed AI-skills gap & 448 & $-0.014\;[-0.096,0.067]$ & $0.119\;[-0.0004,0.238]$ & $0.069\;[-0.033,0.171]$ & $0.065\;[-0.125,0.255]$ \\
\bottomrule
\end{tabular}}
\tabnote{Rows use each alternative's native finite-score support.  The outcome cells, age contrast, calendar, estimator, Webb-software control, and preperiod weighting rule are fixed.  Each beta anchor is rebuilt on exactly the alternative row's support.  Coefficients compare Q5 with Q1; intervals and paired differences use 9,999 common occupation-level Rademacher score multipliers.  A null-containing paired interval is a nondetection, not equivalence.  Continuous-score estimates and raw-to-weighted-SD scaling are preserved in the public result files.}
\end{table}


\begin{table}[!htbp]
\centering
\caption{Post-outcome exploratory direct and software-complement primitives}\label{tab:appendix_primitive_ds}
\small
\setlength{\tabcolsep}{4pt}
\resizebox{\textwidth}{!}{%
\begin{tabular}{llrrr}
\toprule
Units or contrast & Term & Estimate [95\% CI] & Bootstrap SE & $\mathrm{MDE}_{80}$ \\
\midrule
Raw share & $E_1\times$ young $\times$ post & $-0.189\;[-0.354,-0.024]$ & 0.085 & 0.239 \\
Raw share & $E_2\times$ young $\times$ post & $-0.089\;[-0.172,-0.005]$ & 0.043 & 0.121 \\
Preperiod SD & $E_1\times$ young $\times$ post & $-0.031\;[-0.058,-0.004]$ & 0.014 & 0.039 \\
Preperiod SD & $E_2\times$ young $\times$ post & $-0.028\;[-0.054,-0.002]$ & 0.014 & 0.038 \\
One SD in each primitive & $E_1+E_2$ joint change & $-0.059\;[-0.098,-0.020]$ & 0.020 & 0.056 \\
\bottomrule
\end{tabular}}
\tabnote{All entries are post-outcome exploratory and use the same 468 occupations and 9,999 common occupation multipliers.  The joint model includes $E_1$, $E_2$, and standardized Webb software, each interacted with young $\times$ post.  Here $E_1$ is the share of tasks classified as directly accelerable and $E_2$ is the share requiring complementary software; beta satisfies $E_1+0.5E_2$.  The preperiod employment-weighted means (standard deviations) are 0.144 (0.164) for $E_1$ and 0.463 (0.315) for $E_2$.  The common-draw correlation between the two standardized coefficient draws is 0.062.  Component-row MDEs use the analytic occupation-clustered standard errors; the joint-change MDE uses its common-draw standard error.  The final row is the declared one-standard-deviation increase in both primitives, not a paired difference between two technologies.  The parameterization does not measure adoption and is not a causal decomposition.}
\end{table}


\begin{table}[!htbp]
\centering
\caption{Post-outcome exploratory task-weight architectures on common support}\label{tab:appendix_architecture_pairs}
\small
\setlength{\tabcolsep}{4pt}
\resizebox{\textwidth}{!}{%
\begin{tabular}{lrrrr}
\toprule
Architecture pair & Alternative coefficient & Beta coefficient & Alternative $-$ beta [95\% CI] & Paired $\mathrm{MDE}_{80}$ \\
\midrule
$E_1$ versus $E_1+0.5E_2$ & $-0.103$ & $-0.132$ & $0.029\;[-0.044,0.103]$ & 0.106 \\
$E_1+0.25E_2$ versus $E_1+0.5E_2$ & $-0.101$ & $-0.132$ & $0.031\;[-0.005,0.068]$ & 0.055 \\
$E_1+0.75E_2$ versus $E_1+0.5E_2$ & $-0.143$ & $-0.132$ & $-0.011\;[-0.039,0.017]$ & 0.041 \\
$E_1+E_2$ versus $E_1+0.5E_2$ & $-0.156$ & $-0.132$ & $-0.024\;[-0.064,0.017]$ & 0.058 \\
\bottomrule
\end{tabular}}
\tabnote{All entries are post-outcome exploratory Q5--Q1 coefficients estimated on the identical 468-occupation support.  Beta is $E_1+0.5E_2$.  At each alternative weight, normalization, cutoffs, and tie-preserving quintile memberships are recomputed, so a paired movement combines score and membership changes rather than holding the treatment groups fixed.  The intervals use 9,999 common occupation multiplier draws, preserving covariance within each pair.  Every paired interval contains zero; this means that the design does not detect a difference, not that the architectures are economically equivalent.}
\end{table}


\begin{table}[!htbp]
\centering
\caption{Post-outcome exploratory Webb conditioning and support}\label{tab:appendix_webb_support}
\small
\setlength{\tabcolsep}{4pt}
\resizebox{\textwidth}{!}{%
\begin{tabular}{lrrr}
\toprule
Specification & Occupations & Beta coefficient [95\% CI] & Same-support paired change [95\% CI] \\
\midrule
Without Webb, fixed primary support & 468 & $-0.134\;[-0.219,-0.049]$ & $-0.002\;[-0.010,0.007]$ \\
With Webb, fixed primary support & 468 & $-0.132\;[-0.220,-0.044]$ & Reference \\
Without Webb, broader beta-valid support & 490 & $-0.135\;[-0.218,-0.051]$ & Not paired \\
\bottomrule
\end{tabular}}
\tabnote{All entries are post-outcome exploratory.  The first two rows hold occupations, beta cutoffs, and the outcome specification fixed; the displayed paired change is without-Webb minus with-Webb, uses common draws, and has $\mathrm{MDE}_{80}=0.013$.  Its interval containing zero is a nondetection, not equivalence.  The final row relaxes Webb availability and therefore changes support and beta cutoffs; it has no paired interpretation.  Webb conditioning is an occupational-software proxy, not realized software or AI adoption.}
\end{table}


\subsection{Later self-reported use}

An external descriptive check merges the frozen beta score to the author-maintained Real-Time Population Survey occupation workbook.  The workbook pools four waves from August 2025 through May 2026; it is not a preperiod measure or a detailed occupation-by-time panel.  The exact-code merge retains 121 nonsuppressed occupations from the frozen 468, with all five quintiles and 20 major families represented.  The unweighted exposure--adoption correlation is 0.504 and mean adoption differs by 0.267 between Q5 and Q1, although the quintile profile is not monotone because Q4's mean exceeds Q5's.  After major-family means are removed, 26.2 percent of the exposure sum of squares remains and the residualized slope is 0.302.

The workbook supplies pooled point estimates and respondent counts but no detailed-cell standard errors or replicate weights.  Table~\ref{tab:rps_adoption} therefore reports no p-values or confidence intervals, and its respondent-count weighting is only a descriptive sensitivity.  These results show positive alignment between the capability score and later reported use; they do not establish adoption timing, explain the CPS coefficient, or identify an employment effect.

\begin{table}[!htbp]
\centering
\caption{Later self-reported GenAI use by frozen beta exposure}\label{tab:rps_adoption}
\begin{tabular}{lcc}
\toprule
 & Equal occupation & Respondent-count sensitivity \\
\midrule
Pearson correlation & 0.504 & 0.647 \\
Spearman correlation & 0.494 & --- \\
Adoption-on-beta slope & 0.493 & 0.695 \\
Mean adoption, Q1 & 0.263 & 0.250 \\
Mean adoption, Q5 & 0.530 & 0.583 \\
Q5 minus Q1 & 0.267 & 0.333 \\
Within-family exposure SS share & 0.262 & 0.191 \\
Within-family slope & 0.302 & 0.465 \\
\bottomrule
\end{tabular}
\tabnote{The exact-code merge retains 121 nonsuppressed occupations from the frozen 468; all five quintiles and 20 SOC major families are represented.  The RPS workbook pools four waves from August 2025 through May 2026 and suppresses rates based on fewer than 20 respondents.  The second column weights occupations by the workbook's pooled unweighted respondent count, which is neither a survey weight nor a verified inverse-variance weight.  Detailed-cell standard errors and replicate weights are unavailable, so the table reports no inference.  The quantities are descriptive validation of score alignment, not an adoption treatment or an explanation of the CPS coefficient.}
\end{table}


\subsection{What validation would require}

None of the sampling intervals above treats a score as a noisy observation of a validated true exposure.  A useful validation sample would begin at the versioned occupation--task level and oversample both broad families and tasks on which architectures disagree.  It would record separate outcomes for technical capability, realized time savings without complementary software, realized time savings with software, observed workplace adoption, and task importance.  Conflating these labels would reproduce the construct problem the validation is meant to solve.

Labels should come from repeated, blinded, independent expert assessments and, where feasible, audited product-use or task-performance records.  Calibration and evaluation occupations should be separated so that crosswalk and occupation aggregation choices are tested out of sample.  The task description, model snapshot, prompt, complementary tools, rater instructions, occupational bridge, and employment weights all require version records.  Sample size cannot be stated as one universal number: it depends on within-occupation task correlation, rater reliability, unequal task weights, the latent quantity being corrected, and whether the downstream target is a continuous slope or tail assignment.  In the absence of these inputs, the revision specifies the validation design but does not apply a generated-regressor correction.
